\scanpage{4}

\deleted{$\mathcal M$ catégorie abélienne $K$-linéaire, $K\subset\mathbb R$}

\deleted{\unsure{a)}} $\mathcal M$ une $\otimes$-catégorie sur $K\subset\mathbb R$, avec \deleted{\unsure{graduation et}} \inserted{\deleted{\unsure{soit}} un} objet inversible \deleted{$K(1)$ de degré $2$.} $L$

1) Soit $M\in\operatorname{Ob}\mathcal M^{\unsure{i|*}}$, $\varphi:M\otimes M\to L\deleted{\unsure{K(1)}}$. Supposons $\varphi$ symétrique ou alternée, $\varphi$ non dégénérée. Alors $\varphi$ définit sur $\operatorname{End}(M)$ une \uline{anti-involution} $u\mapsto u^*$ $[\,\varphi(ux,y)=\varphi(x,u^*y)\,;$ posant $\varphi(x,y)=\langle x,\widetilde\varphi y\rangle$, \inserted{$\widetilde\varphi:M\to\check M\otimes L$,} on aura donc $u^*=\widetilde\varphi^{-1}\circ\check u\otimes1\circ\widetilde\varphi\,]$. \deleted{\unsure{Tout}} On a une bijection

\[
\operatorname{Hom}(M\otimes M,L\deleted{\unsure{K(1)}})\simeq\operatorname{End}(M)
\]

[$\psi$ et $u$ se correspondent si $\psi(x,y)=\varphi(x,uy)$, i.e $\widetilde\psi=\widetilde\varphi u$], et cette dernière \deleted{\unsure{oblige|bijection}} transforme $\psi\mapsto\varepsilon{}^{\unsure{s|t}}\psi$ en $u\mapsto\varepsilon u^*$, (où $\varepsilon=\pm1$ suivant que $\varphi$ sym. ou alternée). Donc les formes sym (resp. alternées) de $\operatorname{Hom}(M\otimes M,L)$ correspondent aux $u$ tels que $u^*=u$.

2) On dit que $\varphi$ \deleted{\unsure{est p|est une|est définie}} satisfait la \uline{condition de Weil} si \inserted{$\varphi$ (sym. ou \unsure{antisym.|altern.} non dég) est telle que} $\operatorname{Tr}u^*u>0$ pour $u\in\operatorname{End}(M)$, $u\ne0$. Ceci implique, si $\operatorname{End}(M)$ de dim finie sur $K$, que $\operatorname{End}(M)$ est $\frac12$ simple, et pour tt ss-objet $N\hookrightarrow M$, $\varphi|N$ est une forme de Weil.

\inserted{\uline{\unsure{donc que|dire que} $M$ est semi-simple} s'il est de long. finie (tt ss-objet a un complémentaire)}

3) \deleted{\unsure{On}} Soient $\varphi\in\operatorname{Hom}(M\otimes M,L)$, $\psi\in\operatorname{Hom}(N\otimes N,L)$, $\varphi$ et $\psi$ $\varepsilon$-sym. $(\varepsilon=\pm1)$. Alors $\varphi\oplus\psi$ satisfait à condition de Weil ssi $\varphi$ et $\psi$ y satisfont, et si de plus, pour tt $u:M\unsure{-|\to}N$, on a

% slanted
\margin{\unsure{Généraliser|Généralisation} au cas \deleted{\unsure{où les}} \unsure{d'une cat|d'un corps|un nombre} quelconque de facteurs (\unsure{non cyclotomie|en coefficients|ou sommandes})}

\[
\operatorname{Tr}u^*u>0\qquad\text{si }u\ne0
\]

[$u^*$ calculé pour les formes $\varphi$ et $\psi$ sur $M$ et $N$].

% archive sheet number: 3
